\section{Sequential Logic}

\begin{definition}
Sequential logic incorporates storage elements to maintain "state", producing outputs based on both current and past inputs.
\end{definition}

\subsection{Fundamental Blocks}

\begin{definition}
Basic storage element: two inverters connected in a feedback loop.
\end{definition}

\begin{definition}
\textbf{Reset-Set (R-S) latch} (built from cross-coupled NAND gates):
Set ($S=0$): $Q=1$, Reset ($R=0$): $Q=0$.
$S=0, R=0$ is forbidden (metastability).
\end{definition}

\begin{definition}
\textbf{Gated D Latch} adds Write Enable ($WE$). Transparent when $WE=1$ (output $Q$ follows $D$), holds value otherwise.
\end{definition}

\begin{definition}
\textbf{D Flip-Flop} uses two D latches (master-slave) to be \textbf{edge-triggered}. Captures $D$ value exactly at rising edge of clock ($0 \to 1$).
\end{definition}

\subsection{Timing and Finite State Machines}

\begin{definition}
\textbf{Clock (CLK)}: Periodic signal coordinating synchronous state transitions. Period must exceed the worst-case combinational delay.
\end{definition}

\begin{definition}
\textbf{Finite State Machine (FSM)} has:
\begin{itemize}
    \item \textbf{Moore FSM}: Outputs depend \textbf{only} on current state.
    \item \textbf{Mealy FSM}: Outputs depend on current state \textbf{and} current inputs.
\end{itemize}
\end{definition}

\begin{remark}
FSM Implementation:
1. State Register (sequential)
2. Next State Logic (combinational)
3. Output Logic (combinational)
\end{remark}

\begin{example}
State Encodings: \textbf{Binary} (min bits), \textbf{One-Hot} (1 bit/state, simple logic), \textbf{Output} (state bits = output bits).
\end{example}

\begin{definition}
\textbf{Register}: Group of D flip-flops with common clock/enable.
\textbf{Memory}: Storage locations indexed by address. Hardware uses an \textbf{Address Decoder} to select a location and a \textbf{Multiplexer} to route data to the output.
\end{definition}
